\documentclass[12pt]{article}
\usepackage[a4paper,margin=1in]{geometry}\usepackage[T1]{fontenc}
\usepackage{lmodern,microtype}\usepackage{amsmath,amssymb,amsthm,mathtools}
\usepackage{booktabs,array,enumitem,xcolor}\usepackage[numbers,sort&compress]{natbib}
\usepackage[colorlinks=true,linkcolor=blue!55!black,citecolor=blue!55!black,urlcolor=blue!55!black]{hyperref}
\linespread{1.07}
\newtheorem{theorem}{Theorem}[section]\newtheorem{proposition}[theorem]{Proposition}
\theoremstyle{remark}\newtheorem{remark}[theorem]{Remark}
\newcommand{\norm}[1]{\left\lVert#1\right\rVert}

\title{Endpoint Isolation Versus Interlevel Transport Certification\\
\large A Conditioning-Scaled Feasibility Audit for the Edge Quartet}
\author{Liang Wang}\date{July 2026}
\begin{document}\maketitle
\begin{abstract}
Validated Riesz projectors and validated cross-level transport are different
tasks.  This paper separates them before an expensive interval computation.
For a computed simple eigenpair we record the local condition number
$\kappa=\|v\|\|w\|/|w^*v|$, residual $\eta=\|Av-\lambda v\|$, spectral
separation $\delta$, and feasibility ratio
\[
 \beta=2\kappa\eta/\delta.
\]

At each of six physical scale/channel endpoints, every outer-quartet mode
has $\beta<1$ by more than twelve orders of magnitude; the largest endpoint
ratio is $3.25\times10^{-13}$.  This indicates that a dedicated interval
eigenpair or contour validation is numerically plausible.

Using the Haar lift of a coarse eigenvector as an approximate fine
eigenvector gives the opposite result.  Across four adjacent cases the
minimum conditioning-scaled transport ratio is $3.3486$ and the maximum is
$29.4916$; no case has all four ratios below one.  Thus endpoint isolation
should be validated directly, while the naive interlevel homotopy should not
be expected to close.  The ratios are floating feasibility diagnostics, not
interval certificates; no validated projector is claimed.
\end{abstract}

\section{Two certification problems}

The local problem asks whether a quartet at one finite endpoint is isolated
and has the advertised rank.  The transport problem asks whether a
predeclared coarse packet map stays inside the corresponding fine spectral
cluster.  The second requires the first plus an interlevel perturbation
estimate.

RH-203 expresses this distinction exactly through
\begin{equation}\label{eq:riesz-transport}
 P_fJ-JP_c=\frac{1}{2\pi i}\int_\Gamma
 R_f(z)(A_fJ-JA_c)R_c(z)\,dz
\end{equation}
\citep{WangRH203}.  The present paper asks whether the floating scales in
this formula are even compatible with a sharp validated calculation.

\section{Local condition and residual}

For a simple right/left eigenpair of a nonnormal matrix, define
\begin{equation}\label{eq:kappa}
 \kappa(\lambda)=\frac{\norm v_2\norm w_2}{|w^*v|}.
\end{equation}
This quantity is invariant under separate nonzero scalings of $v,w$ and is
the norm of the rank-one spectral projector under biorthogonal
normalization \citep{StewartSun1990,TrefethenEmbree2005}.

For a computed pair $(\widehat\lambda,\widehat v)$ let
\begin{equation}\label{eq:residual}
 \eta=\norm{A\widehat v-\widehat\lambda\widehat v}_2,
 \qquad
 \delta=\min_{\mu\ne\widehat\lambda}|\widehat\lambda-\mu|.
\end{equation}
The dimensionless ratio
\begin{equation}\label{eq:beta}
 \beta=\frac{2\kappa\eta}{\delta}
\end{equation}
compares the conditioning-amplified residual with half the observed
separation.

\begin{remark}[Meaning of $\beta$]
The inequality $\beta<1$ is used here as a feasibility gate, not as a
standalone validated eigenvalue theorem.  A proof still needs outward-rounded
matrix data and an appropriate interval, Krawczyk, Schur, or contour
argument.  Conversely, $\beta\gg1$ warns that this particular residual
budget cannot separate the mode.
\end{remark}

\section{Endpoint result}

The maximum ratio and largest local condition in each endpoint packet are:
\begin{center}
\small
\begin{tabular}{llrr}
\toprule
$\sigma$ & side & max $\beta_{\rm endpoint}$ & max $\kappa$\\
\midrule
$0.04$ & left  & $3.25\times10^{-13}$ & $31.46$\\
$0.04$ & right & $3.11\times10^{-13}$ & $30.69$\\
$0.02$ & left  & $1.09\times10^{-13}$ & $8.39$\\
$0.02$ & right & $7.20\times10^{-14}$ & $8.40$\\
$0.01$ & left  & $1.63\times10^{-13}$ & $17.54$\\
$0.01$ & right & $2.30\times10^{-13}$ & $17.69$\\
\bottomrule
\end{tabular}
\end{center}
All 24 endpoint mode records are deep inside the feasibility gate.  The
projectors are nonnormal but not so ill conditioned that double-precision
eigenpair residuals consume the spectral separation.

This supports a targeted validation strategy: enclose each input matrix,
validate four simple endpoint eigenpairs or one rank-four contour, and then
form outward projector/residue bounds locally.

\section{Lifted coarse modes as fine approximants}

Let $(\lambda_c,v_c)$ be a coarse mode and $J$ the Haar embedding.  Its fine
residual at the unchanged eigenvalue is
\begin{equation}\label{eq:lifted-residual}
 \eta_{c\to f}=\norm{A_fJv_c-\lambda_cJv_c}_2.
\end{equation}
This residual includes both eigenvalue motion and eigenspace rotation.  We
scale it by the condition and separation of the matched fine mode using
\eqref{eq:beta}.

\begin{center}
\begin{tabular}{llrr}
\toprule
step & side & min $\beta_{c\to f}$ & max $\beta_{c\to f}$\\
\midrule
$0.04\to0.02$ & left  & $7.793$ & $20.990$\\
$0.04\to0.02$ & right & $3.349$ & $3.730$\\
$0.02\to0.01$ & left  & $12.015$ & $29.492$\\
$0.02\to0.01$ & right & $5.982$ & $13.856$\\
\bottomrule
\end{tabular}
\end{center}
Every transported mode fails the unit feasibility gate.  This agrees with
the order-one subspace angles of RH-202 and shows that the issue is not
roundoff error in the endpoint eigendecomposition.

\section{Why endpoint validation can succeed while transport fails}

The endpoint residual is generated by solving the fine eigenproblem itself
and is near machine precision.  The transported residual tests a physical
hypothesis about how two different operators are related.  There is no
reason for it to be small unless a renormalization theorem enforces that
relationship.

Thus the logical implication is one-way:
\begin{equation*}
 \text{good interlevel certificate}
 \Longrightarrow \text{isolated endpoints},
\end{equation*}
but isolated endpoints do not imply a good Haar transport certificate.

\section{Avoiding the full eigenvector condition number}

The physical bulk matrices contain removed peripheral directions and highly
clustered interior modes.  The condition number of a complete eigenvector
matrix can be enormous and is not the right local quantity for one isolated
simple mode.  We therefore use the rank-one projector condition
\eqref{eq:kappa}.  A validated contour method would similarly estimate the
resolvent only on the selected contour rather than diagonalize the complete
operator.

\section{Recommended validated workflow}

The finite audit supports the following order:
\begin{enumerate}[leftmargin=2em]
\item enclose the six finite input matrices with outward rounding;
\item validate endpoint quartet counts and separations independently;
\item enclose local Riesz projectors and residues;
\item do not spend interval effort on the unmodified Haar homotopy;
\item first derive a smaller-defect renormalized map or move to a scalar
divisor formulation.
\end{enumerate}
This makes validation proportional to the evidence rather than using an
expensive contour computation to rescue a map already contradicted by
order-one residuals.

\section{Three possible endpoint validators}

There are several compatible rigorous implementations:
\begin{enumerate}[leftmargin=2em]
\item an interval Newton or Krawczyk argument for each simple eigenpair,
including a normalization equation for the eigenvector;
\item a validated Schur decomposition followed by cluster separation and a
Sylvester-equation projector enclosure;
\item an argument-principle or contour-resolvent computation validating the
rank-four cluster as a whole.
\end{enumerate}
The third route avoids labeling individual conjugate modes, while the first
connects most directly to residue enclosures.  A complete implementation
must include rounding error in constructing the physical matrices, not only
the error of the eigensolver applied to already rounded input.

The endpoint ratios suggest that all three methods have numerical room.  The
transport ratios show that none should use the Haar-lifted coarse vector as a
tight initial enclosure for the fine eigenvector.

\section{Claim boundary}

All $\beta$ values use floating eigenvalues, vectors, and separations.  They
are feasibility indicators only.  No interval matrix enclosure, Krawczyk
proof, validated contour rank, or validated Riesz projector appears here.
The endpoint result is positive evidence for such work; the transport result
is a finite negative for the naive map.  Gate A remains open.

\section{Next question}

One possible response is to enlarge the outer quartet into a larger
modulus-selected cloud.  RH-209 tests ranks $2$ through $32$ on every
adjacent case.  If larger clouds absorb the rotating directions, their
principal angles should improve; if not, the selection rule itself must
change.

\bibliographystyle{plainnat}\bibliography{references}\end{document}
